\documentclass[10pt]{article}

% graphicx package, useful for including eps, pdf, and raster graphics
\usepackage{graphicx}
\DeclareGraphicsExtensions{.png,.png,.jpg,.pdf}

% basic packages
\usepackage{fancyhdr}
\usepackage{color}
\usepackage{parskip}
\usepackage{float}
\usepackage{microtype}
\usepackage{url}
\usepackage{hyperref}
\usepackage{booktabs}
\usepackage{makecell}
\usepackage{multirow}
\usepackage{pbox}
\usepackage{enumitem}

% page layout
\usepackage[top=1in, bottom=1in, left=0.75in, right=0.75in, includefoot, heightrounded]{geometry}

% Useful for seeing the layout as you tweak it.
%\usepackage{showframe}

\setlength{\headheight}{1.1in}
\pagestyle{fancy}
\fancyhf{}

% fonts
\usepackage{fontspec}

\usepackage{polyglossia}
\setmainlanguage{russian}
\setotherlanguage{english}
\newfontfamily\cyrillicfont{DroidSans}[
  Extension   = .ttf,
  UprightFont = *,
  BoldFont    = *-Bold ]
\setmainfont{DroidSans}[
  Extension   = .ttf,
  UprightFont = *,
  BoldFont    = *-Bold ]
\setmonofont{DroidSansMono}[
  Extension   = .ttf,
  UprightFont = * ]

\hypersetup{colorlinks,urlcolor=blue}
\urlstyle{same}

\lhead{\pbox[b][1in][c]{\textwidth}{
    \small
    \textbf{Место проведения тестирования:}\\
    Northwest Genomics Center\\
    University of Washington School of Medicine\\
    Genome Sciences Box 355065\\
    3720 15th Ave NE\\
    Seattle, WA 98195-5065}}
\rhead{\includegraphics[height=0.75in,trim=0 -25 0 0]{logo}}

\newcommand{\PageLine}{\rule{\textwidth}{0.25mm}}

\begin{document}
\sloppy % prevent overfilled lines

\begin{description}[font=\normalfont,align=left,labelwidth=12em]
\item [\textbf{\VAR{pat_name|e}}]
\item [\textbf{Дата рождения:}] \VAR{birth_date|e}
\item [\textbf{Образец:}]
  %- if swab_type == "ans"
  Мазок из передный носа,
  %- elif swab_type == "mtb"
  из носа,
  %- else
  \errmessage{Invalid swab_type "\VAR{swab_type}"}
  %- endif

  Самостоятельно взят

\item [\textbf{Штрихкод образца:}] \VAR{qrcode|e}
\item [\textbf{Дата забора образца:}] \VAR{collect_ts|e}
\item [\textbf{Дата отчета:}]
  %- if status_code == "never-tested"
  Обследование не проводилось
  %- else
  \VAR{result_ts|e}
  %- endif
\end{description}

\PageLine

\begin{center}
\Large
\textbf{РЕЗУЛЬТАТ ИССЛЕДОВАНИЯ. Этим результатом стоит поделиться с врачом}
\end{center}

\bigskip

\large \underline{\textbf{Об этом тесте}}

\begin{itemize}
\item

  COVID-19 — это \textbf{заболевание}, вызываемое \textbf{вирусом} SARS-CoV-2.
  Такой метод тестирования направлен на поиск вируса в образце, взятом из вашего
  носа, чтобы определить, есть ли у вас COVID-19. Научное название этого вида
  теста — «полимеразная цепная реакция с обратной транскрипцией» (сокращенно —
  «ОТ-ПЦР»).

\end{itemize}

\bigskip

\large \underline{\textbf{Важно помнить}}

Прежде чем узнавать результат своего теста, вам следует учесть следующие моменты:

\begin{itemize}
\item

  \textbf{Тип образца}, который вы взяли, и \textbf{выполнение теста в
  лаборатории} являются частью \textbf{исследования}, и в этом отчете описан
  \textbf{результат вашего тестирования.}

\item
%- if staff_observed == "f"
  Вы самостоятельно взяли и подписали этот образец (за себя или за кого-то
  другого), а служба доставки передала его в лабораторию. Он не был взят,
  подписан или доставлен медицинским работником (медсестрой или врачом).

  %- elif staff_observed == "t"
  Вы самостоятельно взяли этот образец у себя (или у кого-то другого), а затем
  подписали его и вернули в лабораторию, воспользовавшись услугами курьера или
  другого человека, не являющегося медицинским работником. Он не был взят или
  доставлен медицинским работником (медсестрой или врачом). 

  %- else
  \errmessage{Unknown value: staff_observed = "\VAR{staff_observed}"}

  %- endif

  Это может повлиять на результат вашего теста, и в медицинском учреждении вам
  могут предложить повторное тестирование.

\item

  Участие в этом \textbf{исследовании} не заменяет надлежащее медицинское
  обслуживание. Если вы чувствуете себя очень плохо или беспокоитесь о
  результате теста, вам следует обратиться к врачу или в медицинское учреждение.

\end{itemize}

\bigskip
\PageLine

\large \underline{\textbf{Результат вашего теста на наличие SARS-CoV-2 (COVID-19)}}

%- if status_code == "negative"
Не обнаружен (отрицательный)\\
%- elif status_code == "positive"
Обнаружен (положительный)\\
%- elif status_code == "inconclusive"
Неопределенный\\
%- elif status_code == "never-tested"
Образец для теста забракован\\
%- else
\VAR{0/0} %# poor man's raise
%- endif

\PageLine
\bigskip

\large \underline{\textbf{Что означает результат вашего теста}}

В таблице указаны все варианты результатов, которые может показать тестирование,
и описано, что каждый из них означает. Вы можете распечатать или отправить по
электронной почте \textbf{результат своего теста}, чтобы поделиться с вашим
врачом или медицинским учреждением.

Независимо от результата вашего теста, если вы чувствуете себя плохо, необходимо
оставаться дома и избегать контактов с людьми. \textbf{Если у вас продолжают
проявляться симптомы, вы чувствуете себя очень плохо, обеспокоены по поводу
своего заболевания или у вас есть вопросы о результате вашего теста, следует
обратиться к врачу или в медицинское учреждение.}

\begin{itemize}


\item

  \textbf{Если в результате Вашего теста указано «Не обнаружен
  (отрицательный)»}, это означает, что наша лаборатория не нашла вирус, который
  вызывает COVID-19, в отправленном Вами образце. Ни один тест не является
  идеальным. Существует небольшая вероятность, что Вы могли заразиться вирусом,
  даже если Ваш тест отрицательный.

\item

  \textbf{Если в результате Вашего теста указано «Обнаружен (положительный)»},
  это означает, что наша лаборатория нашла вирус, который вызывает COVID-19, в
  отправленном Вами образце. Вам следует оставаться дома и избегать контактов
  с людьми, особенно если Вы плохо себя чувствуете.

  Результаты вашего теста будут переданы в Washington State Department of Health
  (Департамент здравоохранения штата Вашингтон) и местный департамент
  здравоохранения. Вам \textbf{не нужно} сообщать о своем результате
  самостоятельно. Департамент здравоохранения может связаться с вами, чтобы
  задать вопросы о вашей болезни и о людях, с которыми вы, возможно,
  контактировали, когда были уже больны.

\item

  \textbf{Если в результате Вашего теста указано «Неопределенный»}, это
  означает, что наша лаборатория не уверена в результате Вашего теста. Ни один
  тест не является идеальным. Неопределенный результат получается, когда мы
  обнаруживаем SARS-CoV-2 в Вашем образце, но в недостаточном количестве, чтобы
  утверждать, что у Вас есть COVID-19.

\item

  \textbf{Если в результате Вашего теста указано «Образец для теста
  забракован»}, это означает, что либо в Вашем образце недостаточно материала
  для тестирования, либо произошла ошибка при выполнении теста. Поскольку Ваш
  образец невозможно протестировать, мы не смогли определить, есть ли у Вас
  вирус, вызывающий COVID-19.

\end{itemize}

\bigskip

\large \underline{\textbf{Ресурсы}}

Если вам нужна дополнительная информация о COVID-19, просмотрите перечисленные
ниже ресурсы от партнеров общественного здравоохранения.

\begin{itemize}
\item
  Public Health \textemdash Seattle \& King County (Государственный отдел
  здравоохранения \textemdash Сиэтл и округ Кинг): \\
  \url{https://www.kingcounty.gov/depts/health/covid-19/languages/russian.aspx}

\item

  Washington State Department of Health (Департамент здравоохранения штата Вашингтон):\\
  \url{https://www.doh.wa.gov/Emergencies/Coronavirus/Russian}

\item
  U.S. Centers for Disease Control and Prevention (CDC, Центры по контролю и
  профилактике заболеваний в США):\\
  \url{https://wwwn.cdc.gov/pubs/other-languages?Sort=Lang%3A%3Aasc&Language=Russian}

\end{itemize}

Если вы хотите задать кому-то вопросы по телефону, звоните в одну из
перечисленных ниже служб.

\begin{itemize}

\item

  \textbf{Медицинские вопросы, связанные с COVID-19:} обращайтесь в King County
  Call Center (кол-центр округа Кинг) с 08:00 до 19:00 по номеру
  \textbf{206-477-3977}

\item

  \textbf{Прочие вопросы о COVID-19, в том числе связанные с бизнесом и
  выполнением определенных требований:} обращайтесь в King County COVID-19
  Business and Community Information Line (Информационный центр округа Кинг по
  связанным с COVID-19 вопросам в сферах предпринимательской и общественной
  деятельности) с понедельника по пятницу, с 08:30 до 16:30 по номеру
  \textbf{206-296-1608}

\item

  \textbf{Общие вопросы о COVID-19 в штате Вашингтон:} обращайтесь в Washington
  State Novel Coronavirus Call Center (кол-центр штата Вашингтон по вопросам
  нового коронавируса) с 06:00 до 22:00 по номеру \textbf{1-800-525-0127}

\end{itemize}

\bigskip

\large \underline{\textbf{Методы и ограничения}}

\textbf{Методы}\\
Этот анализ выполняется для обнаружения SARs-CoV-2 с использованием метода ПЦР с
обратной транскрипцией. Такой анализ использует амплификацию (процесс
многократного копирования) вирусных сегментов ORf1B и S, а также выявляет
контроль над амплификацией РНКазы P человека. Разработкой теста и определением
его рабочих характеристик занимался научно-исследовательский центр Northwest
Genomics Center университета University of Washington. Эта лаборатория
лицензирована Медицинский испытательный центр WA Department of Health как
подходящая для проведения клинических лабораторных исследований повышенной
сложности.

\textbf{Ограничения}

\begin{itemize}

\item

  Результаты предоставляются как часть исследования.

\item
%- if staff_observed == "f"
  Вы самостоятельно взяли и подписали этот образец (за себя или за кого-то
  другого), а служба доставки передала его в лабораторию. Он не был взят,
  подписан или доставлен медицинским работником (медсестрой или врачом).

  %- elif staff_observed == "t"
  Вы самостоятельно взяли этот образец у себя (или у кого-то другого), а затем
  подписали его и вернули в лабораторию, воспользовавшись услугами курьера или
  другого человека, не являющегося медицинским работником. Он не был взят или
  доставлен медицинским работником (медсестрой или врачом). 

  %- else
  \errmessage{Unknown value: staff_observed = "\VAR{staff_observed}"}

  %- endif

  Это может повлиять на результат вашего теста, и в медицинском учреждении вам
  могут предложить повторное тестирование.

\item

  Поскольку этот отчет основан на данных научного исследования, ваши результаты
  могут оказаться неточными для диагностирования COVID-19, заболевания,
  вызываемого вирусом SARS-CoV-2.

\end{itemize}

\end{document}
